\documentclass[a4paper,12pt]{article}

\usepackage[spanish]{babel}
\usepackage[utf8]{inputenc}
\usepackage[T1]{fontenc}

\usepackage{geometry}
\geometry{margin=2.5cm}

\usepackage{hyperref}
\usepackage{microtype}

\usepackage{fancyhdr}

\pagestyle{fancy}
\fancyhf{}
\lhead{Class Notes}
\rhead{The Culture of Business}
\cfoot{\thepage}

\title{\vspace{-2cm}\textbf{Class Notes}\\[0.2cm]
\large The Culture of Business}
\author{Juan Ignacio Elosegui}
\date{Universidad Torcuato Di Tella \\ Mayo 2026}

\begin{document}
\maketitle
\section*{Second Part}
\subsection*{Class XIII - 05.05 (Consultants IV: Profitable Uncertainty)}
\noindent \textbf{Reading:} Stein, Chapter 6 -- ``Uncertainty at Work'' \\
\\
\noindent \textbf{Profitable uncertainty} is the idea that consultants benefit from -- and actively cultivate -- uncertainty. It is both a \emph{kind of speech} and a \emph{kind of orientation} toward the future. Uncertainty is not a problem to be solved but a resource to be exploited: it justifies the consultant's presence, generates new engagements, and keeps clients dependent. \\
\\
\noindent Stein identifies \textbf{three mechanisms} through which consultants make uncertainty profitable:
\begin{enumerate}
	\item \textbf{Potential worlds:} Consultants construct possible future scenarios -- some desirable, some threatening -- that make the client feel the need for expert guidance. By painting vivid alternative futures, they create urgency and justify action.
	\item \textbf{Free decisions:} Consultants present \emph{options, not suggestions}. This shifts the burden of decision to the client while allowing the consultant to hedge against failure. If things go wrong, it was the client's choice.
	\item \textbf{Consultant engagement:} The ongoing relationship itself becomes the product. Uncertainty never fully resolves, so the consultant is always needed for the next step.
\end{enumerate}

\noindent \textbf{The PAL (Profit Acquisition Letter)} \\
\noindent Example of a consulting engagement with an African government. The PAL is a document that outlines potential revenue gains the consultant can help capture. It functions as a sales tool: it identifies ``gaps'' -- between current revenue and potential revenue -- and positions the consultant as the one who can close them. \\
\\
\noindent \textbf{``The Gap'' as rhetorical device:} ``Gap'' is widely used in consulting -- even included in company names. It is far more common than ``difference'' as a term. Consultants constantly invoke gaps -- revenue gaps, performance gaps, knowledge gaps. These are illustrated with charts and shapes, comparing an entity to a \emph{potential state}. The gap is a powerful rhetorical tool because it simultaneously diagnoses a problem and implies a solution. Gaps can be ``bridged,'' ``filled,'' or ``closed'' -- and doing so requires the consultant's expertise. What is all of this for? To present to the client a potential world that is \emph{just} out of reach. \\
\\
\noindent \textbf{Example 2: Norwegian management consulting (Henningsen)} \\
\noindent Erik Henningsen's \emph{Management and Morality} -- ethnography of management consultancy seminars in Norway. \\
\\
\noindent \textbf{The ``Over the Edge'' Project in Linnfjell:} A small Norwegian town (Syndinfjeld area) where the local industry (Norindustry) had historically taken care of everything -- ``the fire department, the nursing home... the snow in the streets during the winter.'' When this paternalistic structure collapsed, the community hired consultants (for the third time) to help envision a future.
\begin{itemize}
	\item Six seminars, three hours each; cost: 45,000 USD.
	\item Participants: local government, union reps, high school students, businesspeople, local and national development NGOs.
	\item Goal: describe future scenarios and develop ``proposals for a vision and strategies for the local community.''
\end{itemize}

\noindent \textbf{The consultant Pettersen} asked: what did a thriving community need? Participants answered ``jobs'' or ``night clubs.'' Pettersen redirected: a \emph{positive attitude} was crucial. ``You need faith to make it happen!'' When someone asked whether money (investments) was what they really needed, Pettersen responded: ``don't say money, say \emph{strategy}. Strategy is created where there is energy, and energy is found where there is hope.'' \\
\\
\noindent \textbf{Four future scenarios} were constructed by participants:
\begin{itemize}
	\item \textbf{Alp Village:} full-scale tourism development based on alpine skiing facilities.
	\item \textbf{Neo-industrial Growth:} industrial production revived through R\&D on renewable energy.
	\item \textbf{The Green Valley:} less commercially successful future, shaped by nature conservation, sustainable agriculture, and preservation of cultural heritage.
	\item \textbf{Miserable Linnfjell:} continued depopulation and social decay.
\end{itemize}

\noindent A local businessman, Torstein Moen (Neo-Industrial Growth group), grew so inspired he quit his job that morning. But then the project ended -- the end product was a 77-slide PowerPoint. The final presentation leaned toward ``Alp Village,'' so the ``Green Valley'' group complained. And Moen was now jobless. \\
\\
\noindent Pettersen's defense: the presentation was not what mattered most, but whether the project ``succeeded in opening their eyes to new possibilities, whether it inspired and motivated them and fostered a sense of community among them.'' \\
\\
\noindent \textbf{Connection to profitable uncertainty:} The consulting engagement in Linnfjell illustrates how consultants construct \emph{potential worlds} (four scenarios), offer options rather than solutions, and frame the value of the engagement in terms of ``opening minds'' rather than concrete outcomes. The uncertainty of Linnfjell's future was made profitable -- for the consultant. \\
\\
\noindent \textbf{McKinsey example:} Report \emph{``Automotive revolution -- perspective towards 2030''} uses the same rhetoric: ``Prepare for uncertainty,'' ``A sophisticated level of scenario planning and agility is required.'' Uncertainty is not merely acknowledged but \emph{prescribed} -- it becomes an orientation that demands more consulting. \\

\subsection*{Class XIV - 07.05 (Financial Analysts I: What They Are For and How to Become One)}
\noindent \textbf{Reading:} Stefan Leins, \emph{Stories of Capitalism} -- ethnography of financial analysts at a Swiss Bank in Zurich. \\
\\
\noindent \textbf{Can analysts actually predict the market?} \\
\noindent Leins asks: ``are financial analysts actually able to act in the market in a way that no one else can?'' A long line of research (Working 1934, Kendall 1953, Osborne 1959, Malkiel 1973) has shown that stock prices follow a \textbf{random walk} -- past movements do not predict future ones. Across all researchers that asked this question, the answer was often ``No'': we would be better off flipping a coin than following their advice, since they cannot consistently outperform the market.
\begin{itemize}
	\item In 1999, a monkey literally outperformed 6,000 professional financial analysts.
	\item 2003--2009: the \emph{Chicago Sun} ran an experiment publishing forecasts from ``Adam Monk'' (actually a monkey) under the title of ``financial expert.'' He outperformed the indexes four years in a row.
	\item Two monkeys and even a cat named Orlando beat the markets.
\end{itemize}

\noindent Leins: ``\emph{If these people have empirically proved to be wrong, why do they even exist?}'' Why do we keep paying them? \\
\\
\noindent \textbf{So what are analysts actually for?} \\
\noindent Analysts are \emph{not there to be ``right.''} They are there to \textbf{make things make sense}. One thing is to be right, and another thing is to make sense -- the stories they create make sense. Their role is to construct plausible narratives that clients and colleagues can act upon. \\
\\
\noindent \textbf{Fundamental analysis:} One of many ways of doing financial analysis -- looking at a financial product (stocks, etc.) and saying what it is really worth.
\begin{itemize}
	\item Based on the distinction between \textbf{intrinsic value} (the value an asset has inherently) and \textbf{market value} (the value at a given point in the market).
	\item If intrinsic value > market value $\rightarrow$ buy (undervalued). If intrinsic value < market value $\rightarrow$ sell (overvalued).
	\item To do this, analysts need data or must estimate it. The key question: \emph{what} is the difference between intrinsic and market value, and \emph{where} does it come from? They need to \textbf{put into words this explanation} -- verbalizing reasons that don't come from equations or models, but from gut feeling and intuition.
	\item Multiple valuation formulas exist (DCF, P/E ratios, etc.), and there is no objectively ``better'' one. Analysts choose based on \textbf{``market feeling''} -- an intuitive, experience-based sense of which model fits the current moment.
	\item Valuing things in finance takes individual skill (art) and chance (luck). \textbf{Numbers don't tell stories.}
\end{itemize}

\noindent \textbf{Narrative precedes calculation:} \\
\noindent According to Copeland et al., the analyst first builds a \emph{story} about what is happening and why, and only then selects numbers and models to support it. The narrative comes before the math. The spreadsheet is a tool of persuasion, not of discovery. \\
\\
\noindent \textbf{Recruitment and rites of passage:} \\
\noindent Preconditions to become a financial analyst: ``be able to work 20-hour shifts for three days in a row, how to calculate a certain cash metric, speak less, make a solution out of thin air, have to be handy with an Excel spreadsheet.'' \\
\\
\noindent In Switzerland, a network of relevant universities exports professionals to the banking industry. All analysts have a master's in finance. Leins: ``For banking, banks go to the source'' -- the only two universities that matter are UZH and St. Gallen. Banking networks are more refined than the Jockey Club in Argentina. \\
\\
\noindent \textbf{Marcel's recruitment process} (how a senior analyst identifies future analysts):
\begin{enumerate}
	\item At networking events, Marcel looks for the student willing to \textbf{summarize the event} at the end of the day. That person -- able to keep a story together -- enters Marcel's focus group.
	\item Selected students are offered \textbf{summer internships} to test whether they can adopt the analyst lifestyle -- not the workload itself, but commitment to the future job.
	\item After the internship, they are sent back to their universities to \textbf{earn the master's degree}.
	\item Once graduated, candidates apply and are \textbf{assessed online} for interpersonal, numerical, and linguistic skills. Very narrow margin to pass.
	\item The \textbf{first interview} asks questions like ``What is the GDP of Ghana?'' (no one knows). The focus is not academic achievement but \textbf{rhetorical skills, imagination, and self-confidence} -- the ability to build a realistic narrative on the spot.
	\item Once offered a position, junior staff must pass \textbf{three examinations}, each requiring memorization of three to five finance textbooks. Long working hours and self-discipline -- a \emph{rite of passage}.
\end{enumerate}

\noindent They look beyond grades: they want rhetorically strong, self-confident, competitive, and \emph{discrete} people. Analysts are trained to be gentlemen -- a very different job from a consultant. \\
\\
\noindent The \textbf{CFA (Chartered Financial Analyst)} credential serves as further symbolic capital, signaling expertise and discipline.

\noindent \textbf{Symbolic capital and education:} \\
\noindent No other Swiss Bank department has as many highly educated people as financial analysis. This is partly due to the complexity of the work, but also because analysts need to \textbf{present themselves as the ones who are able to understand the market}. A master's degree, PhD, or CFA underlines their symbolic capital. Education also grants access to \textbf{social networks} -- this is a very network-oriented job. \\

\subsection*{Class XV - 12.05 (Financial Analysts II: The Stories that Hold the Economy Together)}
\noindent \textbf{Reading:} Leins, \emph{Stories of Capitalism} (continued). \\
\\
\noindent The job of a financial analyst is to \textbf{tell good stories}. \\
\\
\noindent \textbf{The search for information:} \\
\noindent Financial analysts rely on a wide range of sources:
\begin{itemize}
	\item \textbf{Newspapers and news services} (Bloomberg, Reuters) -- the baseline that everyone reads.
	\item \textbf{Company reports} -- structured documents that present a company's financial health, strategy, and outlook.
	\item \textbf{Personal networks} -- contacts cultivated at conferences, through business cards, and via alumni connections. These provide informal, sometimes privileged, information.
\end{itemize}

\noindent \textbf{The temporality of information:} \\
\noindent Information has a \textbf{shelf life}. Once information is widely known, it is ``priced in'' -- already reflected in market prices -- and therefore no longer useful for gaining an edge. The value of information decays rapidly. Analysts must constantly search for information that is \emph{not yet} priced in. \\
\\
\noindent \textbf{Examples of information temporality:}
\begin{itemize}
	\item \textbf{Hollande's election (France):} Markets reacted before analysts could fully process the implications. By the time an analyst writes a report, the market has already moved.
	\item \textbf{Fukushima disaster:} The nuclear catastrophe instantly affected energy stocks worldwide. Speed of reaction was everything.
	\item \textbf{Arab Spring:} Political upheaval created uncertainty in markets. Those who anticipated it profited; those who reacted after the fact were too late.
\end{itemize}

\noindent \textbf{The tension of originality:} \\
\noindent Analysts face a paradox: they all read the same newspapers, use the same Bloomberg terminals, and attend the same conferences. Using \textbf{common sources reduces risk} (if everyone agrees, you are less likely to be singled out for a bad call), but to add value, an analyst must \textbf{stand out} -- offer an interpretation that others have not seen. Originality is both demanded and dangerous. \\
\\
\noindent \textbf{Company reports:} \\
\noindent Company reports are a key genre. They follow a standardized structure: financial statements, management discussion, outlook. But analysts do not simply read them -- they \emph{interpret} them, looking for signals, omissions, and tone. The report is not transparent data; it is a \textbf{crafted narrative} that the company wants to tell about itself, and the analyst's job is to read between the lines. \\
\\
\noindent \textbf{Stories holding the economy together:} \\
\noindent Leins's central argument: the economy does not run on numbers alone. It runs on \textbf{stories} -- narratives that make sense of data, that connect past events to future expectations, that give actors reasons to buy, sell, or hold. Financial analysts are, fundamentally, \emph{storytellers}. Their stories do not merely describe the economy -- they \textbf{constitute} it, shaping expectations and actions that move markets. \\

\end{document}
